% Generated by report_gen.py (ai-ee v1) - do not hand-edit.
\documentclass[11pt,a4paper]{article}
\usepackage[margin=2.2cm]{geometry}
\usepackage{graphicx}
\usepackage{booktabs}
\usepackage{longtable}
\usepackage{pdfpages}
\makeatletter
\@ifundefined{KV@Gin@artifact}{\define@key{Gin}{artifact}[]{}}{}
\makeatother
\usepackage{xcolor}
\usepackage[hidelinks]{hyperref}
\setcounter{tocdepth}{1}
\begin{document}
\sloppy
\begin{center}
{\LARGE\bfseries bb-ldo --- Design Document}\\[6pt]
{\large ai-ee v1 pipeline}\\[2pt]
generated 2026-08-16 23:29:34
\end{center}
\tableofcontents

\section{Board and Run Metadata}
{\small
\begin{longtable}{lp{11cm}}
\toprule
\textbf{Field} & \textbf{Value} \\
\midrule
\endhead
board & bb-ldo \\
workspace & \texttt{boards/bb-ldo} \\
phase & P4 \\
gate status & all 1 recorded gates pass \\
state created & 2026-08-16T17:24:47 \\
state updated & 2026-08-16T23:29:14 \\
\bottomrule
\end{longtable}
}
\section{Overview}
\subsection*{bb-ldo brief (owner, verbatim)}
learning block-basics: a linear regulator. Input 5 V DC on a screw terminal; output 3.3 V at up to 500 mA on a screw terminal. Must hold regulation at full load with no airflow. Workspace boards/bb-ldo.

\subsubsection*{Resolved mode (state.py mode, P0)}
\begin{itemize}
\item token: \texttt{learning block-basics:}
\item target: block-basics
\item scope tier: block-only (excludes protection, filtering, indicators,
test-points, config, second-rail, mechanical, enclosure-fit)
\item binding: canonical -\textgreater{} geometry is an OUTPUT (board\_init --outline auto,
place, then \texttt{board\_edit --outline fit})
\item stage under study: none
\end{itemize}
\subsection*{Stackup}
Chosen stackup: \texttt{JLC2313\_1.6} - 2 layers, 1.6 mm, 1 oz outer copper, HASL
\section{Requirements}
\subsection*{bb-ldo - requirements (P0)}
Source: \texttt{brief/brief.md} (owner, verbatim). Mode resolved at P0 and recorded in \texttt{state.json}. Contract: \texttt{reference/build-modes.md}.

\subsubsection*{1. Function}
A single-block linear regulator bench board: it takes 5 V DC in on a screw terminal and delivers a regulated 3.3 V rail at up to 500 mA out on a screw terminal, holding regulation at full load in still air (no airflow, no heatsink other than the board's own copper).

\textbf{Build mode.} Token \texttt{learning block-basics:} -\textgreater{} target \textbf{block-basics}, scope tier \textbf{block-only}, binding \textbf{canonical}, stage under study: \textbf{none} (the run teaches the block end to end, not one stage).

\begin{itemize}
\item Scope tier \texttt{block-only}: the board is the regulator, one input interface and
one output interface, plus exactly the support parts the regulator's
datasheet requires at this operating point. Excluded feature classes -
protection, filtering (beyond what the datasheet requires), indicators,
test-points, config, second-rail, mechanical, enclosure-fit - are SCOPE
decisions, not engineering judgements, and their absence is not a finding.
\texttt{thermal} is excluded by no tier: see sections 4 and 5.
\item Binding \texttt{canonical}: the board geometry is an \textbf{OUTPUT} of the design, not
an input. P5 runs \texttt{board\_init --outline auto}, P6 places to the canonical
thermal/loop layout, then \texttt{board\_edit --outline fit} makes the board what
the placement earned. No dimension in this document is a HARD cap.
\item The mode relaxes geometry, cost and packaging only. Every electrical spec,
gate, coverage check, research trigger, DFM check and safety question below
is at full rigor.
\end{itemize}
\subsubsection*{2. Interfaces}
\begin{quote}\small\ttfamily\raggedright
\textbar{} \# \textbar{} interface \textbar{} direction \textbar{} electrical \textbar{} connector \textbar{}\\
\textbar{}---\textbar{}---\textbar{}---\textbar{}---\textbar{}---\textbar{}\\
\textbar{} J-in \textbar{} DC input \textbar{} in \textbar{} 5 V DC nominal, \textasciitilde{}505 mA max (see section 3) \textbar{} screw terminal, 2 position (+5V, GND) - STATED by owner \textbar{}\\
\textbar{} J-out \textbar{} 3.3 V rail \textbar{} out \textbar{} 3.3 V DC, 0 to 500 mA \textbar{} screw terminal, 2 position (+3V3, GND) - STATED by owner \textbar{}
\end{quote}
\begin{itemize}
\item No other external interface: no USB, no RF, no signal I/O, no enable/config
pin brought out, no status LED (indicators are excluded at \texttt{block-only}; a
regulator enable/config strap is \texttt{config}, also excluded - the block runs
always-on with EN tied to its always-on state as the datasheet requires).
\item No test points beyond the two terminals: output voltage and load current are
measurable at J-out, which satisfies the block's own measurement need.
\item Screw terminal pitch, wire range and mounting technology are unstated -
question 8. ASSUMED until answered: 2 position, 5.08 mm pitch, through-hole,
wire 0.2-2.5 mm2, from the JLCPCB catalog.
\end{itemize}
\subsubsection*{3. Power}
\begin{itemize}
\item \textbf{Input}: 5 V DC, single source, on J-in. Source identity and fault current
are UNSTATED and safety-relevant - questions 1 and 2. Voltage tolerance is
UNSTATED - question 4. ASSUMED until answered: 4.75-5.25 V (5 V +/- 5\%),
which sets the worst-case dropout headroom at 1.45 V at 500 mA.
\item \textbf{Output rail (the only rail)}: 3.3 V, 0 to 500 mA. Accuracy UNSTATED -
question 5. Continuous vs peak current UNSTATED - question 6. Load character
and any load-step/transient requirement UNSTATED - question 7.
\item \textbf{Derived budget (GUESS - derived arithmetic, not owner-stated)}:
\begin{itemize}
\item Output power at full load: 3.3 V {*} 0.5 A = 1.65 W.
\item Regulator dissipation, worst case high line: (5.25 - 3.3) {*} 0.5 = \textbf{0.98 W};
nominal line: (5.0 - 3.3) {*} 0.5 = \textbf{0.85 W}. Plus quiescent, so call it
\textasciitilde{}1 W to be dissipated by the package and the board copper.
\item Input current at full load: \textasciitilde{}505 mA (load + quiescent). Efficiency is
inherently \textasciitilde{}63-66\% - a property of the topology the owner chose, not a
requirement to meet.
\end{itemize}
\item No second rail, no sequencing, no soft-start requirement, no power-good.
\item No battery, no charging - subject to question 1.
\end{itemize}
\subsubsection*{4. Environment}
\begin{itemize}
\item \textbf{Still air, no airflow} - STATED by owner ("must hold regulation at full
load with no airflow"). No forced convection, no heatsink, no thermal
interface to anything but the PCB.
\item Ambient \textbf{0 to 50 C, indoor bench, no enclosure} - mode default for
\texttt{block-only} (taken, not asked). The thermal design point is therefore
\textbf{\textasciitilde{}1 W dissipated at 50 C ambient in still air}.
\item No ingress rating, no vibration, no shock, no conformal coating, no altitude
or humidity spec: bench use.
\item Thermal is in scope at every tier. The regulator must stay within its own
safe junction temperature at the design point above, with margin - i.e. the
required theta\_JA follows from the package's Tj(max) and that 1 W at 50 C
(for a 125 C part this is roughly 75 C/W or better). Meeting it is a P2/P3
decision (package choice, copper area, via stitching); it is a REQUIREMENT
here, not an optimization.
\end{itemize}
\subsubsection*{5. Size \& mounting}
\begin{itemize}
\item {*}{*}No HARD size cap. Board outline is RELAXABLE (canonical) - it is an OUTPUT
of the design, not an input.{*}{*} The owner stated no dimension, and under the
\texttt{canonical} binding none would bind anyway: P5 opens with
\texttt{board\_init --outline auto} (generous provisional room), P6 places the
canonical layout, \texttt{board\_edit --outline fit --margin M} then makes the board
what that placement earned, and P7 routes it. Re-run \texttt{planes\_gen} if the
board grew.
\item The outline is the radiator. The copper area needed to hit the theta\_JA in
section 4 is a primary input to what the placement earns - it is exactly the
mechanism that must not be pre-empted by a guessed number (the bb-buck
defect: a stated size that bound, and placement optimizing to fit it).
\item Height: unconstrained (no enclosure). The tallest part will be the screw
terminals.
\item Mounting: mechanical features beyond mounting the bare board are excluded at
\texttt{block-only}. Mounting holes are UNSTATED - question 9. ASSUMED until
answered: none.
\item Layer count: mode default, "the fewest layers the block honestly needs".
ASSUMED 2 layers, with thermal copper on both and stitching vias under the
regulator's pad; confirmed or revised at P2 against the thermal requirement.
\end{itemize}
\subsubsection*{6. Quantity \& budget}
\begin{itemize}
\item Quantity \textbf{5} (mode default, taken not asked).
\item Cost: \textbf{minimal} (mode default). No unit-cost target; the mode relaxes cost.
\item Stop at \textbf{P9} (fab package + DFM). Ordering is a separate owner decision.
\end{itemize}
\subsubsection*{7. Assembly}
\begin{itemize}
\item \textbf{JLCPCB PCBA, single-sided assembly} (mode default, taken not asked). All
SMT parts on the top side.
\item ASSUMED: if the chosen screw terminal is through-hole only, it is hand
soldered after PCBA (JLC economy assembly is SMT-only) - a packaging
detail the mode relaxes. An SMT-mount terminal in the JLC catalog is
preferred if one exists at the required current and wire size.
\item Standard JLCPCB 2-layer process, HASL or ENIG per default DFM, no special
stackup, no controlled impedance, no thick copper unless P2's thermal work
shows it is required (in which case it is a datasheet-driven requirement,
not a feature).
\end{itemize}
\subsubsection*{8. Compliance/safety flags}
Assessed against the brief; every flag here is at full rigor - no mode relaxes any of it.

\begin{quote}\small\ttfamily\raggedright
\textbar{} flag \textbar{} applies? \textbar{} basis \textbar{}\\
\textbar{}---\textbar{}---\textbar{}---\textbar{}\\
\textbar{} Mains voltage \textbar{} {*}{*}NOT APPLICABLE as stated{*}{*} - but the SOURCE of the 5 V is unstated (question 1). If it is a mains-derived adapter, isolation lives in the adapter and the board stays SELV. \textbar{} brief says 5 V DC in \textbar{}\\
\textbar{} Battery / charging \textbar{} {*}{*}UNKNOWN - must be answered before P2{*}{*} (question 1). No chemistry, charging or protection may be guessed. \textbar{} source unstated \textbar{}\\
\textbar{} \textgreater{} 30 V anywhere \textbar{} No. Max node is the 5 V input. \textbar{} brief \textbar{}\\
\textbar{} High current (\textgreater{} 3 A) \textbar{} No. 500 mA out, \textasciitilde{}505 mA in. \textbar{} brief \textbar{}\\
\textbar{} Motors / inductive loads \textbar{} UNKNOWN - depends on the load (question 7). A motor or solenoid load changes the transient and reverse-EMF picture. \textbar{} load unstated \textbar{}\\
\textbar{} RF transmit \textbar{} No. \textbar{} brief \textbar{}\\
\textbar{} Hot surfaces \textbar{} {*}{*}YES, advisory.{*}{*} \textasciitilde{}1 W in a small package in still air: the regulator and its copper will be hot to the touch (typically 80-110 C surface at the design point). Bench handling hazard, recorded so the owner and the reviewers see it; not a scope item to fix. \textbar{} derived, section 3 \textbar{}\\
\textbar{} Fault behaviour \textbar{} The `block-only` tier excludes reverse-polarity, over-voltage and fusing protection on the screw terminal. A screw terminal invites a miswire, and without protection a reversed or over-voltage input destroys the block and can dump the source's full fault current. This is a recorded SCOPE consequence - question 3 asks the owner to accept it explicitly, and questions 1-2 bound how bad the fault can get. \textbar{} tier + brief \textbar{}
\end{quote}
The pipeline does not proceed past P2 on guessed answers to questions 1-3.

\subsubsection*{9. Open questions}
Safety - must be answered (no default is taken without your word):

1. \textbf{What supplies the 5 V?} (a) bench power supply, (b) USB wall adapter or USB port, (c) a battery pack, (d) something else. {*}Default if you have no preference: (a) a bench supply.{*} If the answer is (c), we also need the battery chemistry, cell count and whether this board ever charges it - none of that can be assumed. 2. \textbf{How much current can that source push into a dead short?} i.e. what is its current limit or fuse rating. {*}Default: 2 A or less (typical bench supply limit / USB adapter).{*} 3. \textbf{Do you accept no input protection?} The block-only scope means no reverse-polarity, over-voltage or fuse protection on the screw terminal: if the input is wired backwards or fed the wrong adapter, the board dies (and the source dumps its fault current into it). {*}Default: yes, accepted - it is a bench block.{*} Say no and we add protection, which moves the board off the block-only scope.

Electrical - defaults offered, correct me if any is wrong:

4. \textbf{What is the input voltage range?} \emph{Default: 4.75 to 5.25 V (5 V +/- 5\%).} This matters: at 4.75 V there is only 1.45 V of headroom above 3.3 V, so a wider low end forces a low-dropout part; at 5.25 V the regulator gets hotter. Tell me the real minimum and maximum if the supply is loose. 5. \textbf{How accurate does the 3.3 V need to be?} {*}Default: +/-3\% (3.20 to 3.40 V) including line, load and temperature drift.{*} 6. \textbf{Is 500 mA continuous, or a short peak?} {*}Default: continuous, 100\% duty, forever - the thermal design is sized for that.{*} If it is a brief peak with a low average, say so; the board can be smaller. 7. \textbf{What is on the output?} {*}Default: a static resistive bench load - no load-step or transient-response requirement, and only the output capacitor the regulator's datasheet requires for stability.{*} If the load switches hard, or is a motor/solenoid, say so - that is a different set of requirements.

Packaging - answer only if you care, defaults apply otherwise:

8. \textbf{Screw terminal preference?} {*}Default: 2 position, 5.08 mm pitch, through-hole, accepting 0.2-2.5 mm2 wire, from the JLCPCB catalog (hand soldered after assembly).{*} Name a pitch, wire size or specific part if you have one. 9. \textbf{Mounting holes?} \emph{Default: none - it is a bare bench block.} Say "two M3" or "four M3" if you want to bolt it down.

\paragraph*{Answers - owner, 2026-08-16 (P0 close)}
All nine answered; each is recorded as a \texttt{state.py decision}. Nothing below is a guess: questions 1-3 (safety) carry the owner's own word.

\begin{quote}\small\ttfamily\raggedright
\textbar{} \# \textbar{} answer \textbar{} consequence \textbar{}\\
\textbar{}---\textbar{}---\textbar{}---\textbar{}\\
\textbar{} 1 \textbar{} {*}{*}Current-limited bench supply.{*}{*} Not a battery, not mains-direct. \textbar{} Section 8 battery flag closes NOT APPLICABLE on an owner answer. No chemistry/charging requirements enter the design. \textbar{}\\
\textbar{} 2 \textbar{} Source fault current {*}{*}\textless{}= 2 A{*}{*}. \textbar{} Bounds the miswire/short energy; still no fusing at this tier (question 3). \textbar{}\\
\textbar{} 3 \textbar{} {*}{*}No input protection accepted.{*}{*} \textbar{} Tier exclusion stands. Reviewers must not report absent protection as a finding. Advisory hazard recorded. \textbar{}\\
\textbar{} 4 \textbar{} {*}{*}4.75-5.25 V{*}{*} (5 V +/- 5\%). \textbar{} 1.45 V worst-case headroom at 500 mA -\textgreater{} a standard (non-LDO) regulator is admissible. Worst-case dissipation {*}{*}0.975 W{*}{*} at high line. \textbar{}\\
\textbar{} 5 \textbar{} {*}{*}+/-3\%{*}{*} (3.20-3.40 V) over line, load, temperature. \textbar{} Ordinary fixed-output regulator accuracy; no trim, no external divider needed if a fixed 3.3 V part is chosen. \textbar{}\\
\textbar{} 6 \textbar{} {*}{*}Continuous, 100\% duty.{*}{*} \textbar{} The thermal design is sized for steady state, not bursts. This is what earns the outline at P6. \textbar{}\\
\textbar{} 7 \textbar{} {*}{*}Static resistive bench load.{*}{*} No transient spec. \textbar{} Only the datasheet's own stability capacitance. No bulk output capacitor is added for transients this board will never see. No motor/inductive load - section 8's motor flag closes NOT APPLICABLE. \textbar{}\\
\textbar{} 8 \textbar{} Default taken: 2 position, {*}{*}5.08 mm pitch, through-hole{*}{*}, 0.2-2.5 mm2 wire, JLCPCB catalog. \textbar{} Hand soldered after PCBA if the chosen part is THT-only. \textbar{}\\
\textbar{} 9 \textbar{} {*}{*}No mounting holes.{*}{*} \textbar{} Bare bench block; mechanical stays excluded. \textbar{}
\end{quote}
\textbf{Owner delegation (2026-08-16):} "Make decisions yourself based on what will give the best learning for the basic system." Design judgment calls and the H1-H4 checkpoint ceremony are delegated to the orchestrator, which presents digests with its rulings instead of blocking questions. NOT delegated and unchanged: every gate, the P2/P3 coverage checks and the research they trigger, DFM, and any NEW safety question that arises.

\textbf{Design point, frozen here:} 0.975 W dissipated at Ta = 50 C in still air, no heatsink, no airflow, continuous. The required theta\_JA and the copper area that buys it are the primary drivers of what the placement earns.

\section{Architecture}
\subsection*{Decisions of record (state.json)}
{\small
\begin{longtable}{lp{5.2cm}p{7.2cm}}
\toprule
\textbf{Phase} & \textbf{What} & \textbf{Why} \\
\midrule
\endhead
P0 & Build mode: mode block-basics; scope block-only; binding canonical & token 'learning block-basics:' in the brief; contract in reference/build-modes.md \\
P0 & Owner delegation: orchestrator rules on design judgment calls, targeting best learning value for the basic block & Owner instruction 2026-08-16: 'Make decisions yourself based on what will give the best learning for the basic system'. Scope of the delegation: engineering judgment and checkpoint ceremony (H1-H4 presented as digests with rulings, not blocking questions). NOT delegated: safety answers (owner-answered below, never guessed), gates, coverage, research, DFM. \\
P0 & Input source: current-limited bench supply, 5.00 V nominal, range 4.75-5.25 V (5 V +/-5\%); source current limit assumed \textless{}=2 A & Owner-answered at the P0 question batch (safety section 8): NOT a battery, no charging, no mains. Section 8 is therefore clear on an owner answer, not a guess. 4.75 V worst case leaves 1.45 V of headroom at 500 mA, so a standard (non-LDO) regulator is admissible; 5.25 V worst case sets dissipation at 0.975 W. \\
P0 & Load: 500 mA CONTINUOUS at 100\% duty into a static resistive bench load; no transient/load-step spec; output tolerance +/-3\% (3.20-3.40 V) over line, load and temperature & Owner-answered. Continuous full load in still air is the brief's own stress case ('must hold regulation at full load with no airflow'), so the thermal design is sized for steady-state \textasciitilde{}1 W, not for bursts. A static load means only the regulator datasheet's own stability capacitance is required - no bulk output capacitor is added for transients the board will never see (block-only tier: nothing the datasheet does not require). \\
P0 & Thermal design point: 0.975 W dissipation at Ta = 50 C, still air, no airflow, no heatsink & Worst-case dissipation = (5.25 - 3.3) V {*} 0.5 A = 0.975 W plus quiescent. Ambient 50 C is the block-only mode default (bench, indoor 0-50 C) and thermal is in scope at EVERY tier - it is support the datasheet requires at the stated operating point, never a feature. This number, not a stated dimension, is what earns the board outline at P6 (binding = canonical). \\
P0 & No input protection: no reverse-polarity device, no OVP, no fuse on the input screw terminal. Advisory hazard: a backwards or wrong-adapter input destroys the board, and the regulator surface runs hot (est. 80-110 C) at the design point & block-only scope tier excludes protection of every kind; this is a SCOPE decision for a bench study board, never an engineering claim that protection is unnecessary in general (build-modes.md knowledge guardrail). Recorded so the exclusion is visible at H1 and so no reviewer reports absent protection as a finding. \\
P0 & Mode defaults taken without asking: qty 5, minimal cost, JLC PCBA single-sided assembly, bench 0-50 C indoor, fewest honest layers (2 expected), no mounting holes, no enclosure, stop at P9. Interfaces: 2x 2-position 5.08 mm through-hole screw terminal (in, out) & block-only tier defaults are apply-do-not-ask (build-modes.md). Mounting holes and enclosure fit are excluded mechanical classes. Screw terminals are the brief's stated interface; 5.08 mm THT is the simplest catalog part that carries 500 mA and accepts bench wire. \\
P1 & Lead part: AMS1117-3.3, SOT-223, LCSC C6186 (JLC Basic, \$0.20 @qty5, stock 1.36M). Runner-up MCP1825S-3302E/DB (C148031) rejected. & Chosen for BOARD-CLASS-APPLICABLE thermal data, which is the mechanism this run exists to teach. AMS1117's datasheet Table 1 is a copper-area -\textgreater{} theta\_JA curve measured on 1/16in FR-4 with 1 oz copper (55 C/W at 2500mm2 top+back, 60-65 C/W at \textasciitilde{}1000mm2, 90 C/W at minimum pad) - i.e. OUR board class, so the outline can be derived from it. MCP1825S is the better part on paper (dropout specified at exactly 500 mA, ceramic-stable, +/-2.5\%), but its only theta\_JA figure (62 C/W) is measured on a JEDEC JESD51-7 FOUR-LAYER board with internal planes; on a 2-layer bench board that number is not achievable and the datasheet gives no curve to design against. Choosing it would mean sizing the copper with no applicable data - exactly the error this board is built to avoid. Secondary: Basic tier (no feeder fee), 14x cheaper, and its two 'weaknesses' are high-value lessons (tab = VOUT forces an isolated heatsink pour; ESR-dependent compensation forces a min-ESR/tantalum output cap instead of an all-ceramic BOM). \\
P1 & Thermal design target: theta\_JA \textless{}= 66 C/W, i.e. Tj \textless{}= 115 C at Pd 0.975 W and Ta 50 C (10 C margin to the AMS1117 125 C steady-state limit). Per datasheet Table 1 this needs roughly 1000 mm2 of top-side copper plus backside copper - the copper area is the primary driver of the outline the placement must earn. & Tj(max) 125 C steady state; running at the rating is not meeting a requirement. At the headline 90 C/W (minimum pad) Tj = 138 C and the part FAILS outright - so the copper area is not decoration, it is the design. This number, not any stated dimension, sizes the board (binding = canonical, geometry is an OUTPUT). \\
P1 & Output tolerance: the board's spec is the part's guaranteed over-temperature window, 3.201-3.399 V (+/-3.0\%, AMS1117 boldface limits). Strict arithmetic stacking of line (10 mV max) and load (25 mV max) regulation on top reaches \textasciitilde{}3.166-3.434 V (+/-4.1\%) - accepted, and visible at H1. & The +/-3\% figure was MY offered P0 default, not an owner-originated spec; the brief itself says only '3.3 V at up to 500 mA'. AMS1117 meets +/-3\% as the datasheet specifies Vout over temperature, and exceeds it only under worst-case arithmetic summation of regulation terms that the Vout limits largely already bound. For a bench block into a static resistive load this is immaterial. Recorded rather than silently re-specified: this is an owner-delegated engineering call, NOT a mode relaxation (a mode relaxes geometry, cost and packaging only). \\
P1 & Output capacitor must be min-ESR (solid tantalum class, 22 uF per datasheet), NOT an all-ceramic MLCC. Carried as a hard BOM constraint into P3/P4. & AMS1117 is a bipolar-NPN pass device compensated for a specific ESR window; its datasheet states a 22 uF solid tantalum output capacitor ensures stability over all operating conditions. A pure low-ESR ceramic output cap is a known instability mode for this family. This is a datasheet REQUIREMENT at the stated operating point, so it is in scope at block-only (support the datasheet requires, not added filtering). \\
P1 & Keep the honest thermal target dt\_c = 65 C (theta\_JA 65 C/W at Pd 1.00 W) even though check\_thermal CANNOT pass it on 2 layers. Plan: accept the P8 thermal\_area finding and waive it with the two vendor tables as evidence; do NOT relax dt\_c to 75 to make the gate green. & VERIFIED against the checker source, not taken on report: check\_thermal models theta\_JA(A) = 55 + (174-55){*}exp(-A/350) for 1oz/2-layer and CLAMPS credited area at A\_SAT = 645 mm2, so its 2L floor is 73.85 C/W - it stops crediting copper exactly where AMS1117 Table 1 keeps improving (1000 mm2 -\textgreater{} 65 C/W, 2500 mm2 -\textgreater{} 55 C/W) and TI's LM1117 SOT-223 1oz top-only sweep agrees (645 mm2 -\textgreater{} 66 C/W). The model is \textasciitilde{}12\% pessimistic beyond its clamp; that is a MODEL limitation, not a board property. Relaxing dt\_c to 75 would state Tj \textless{}= 125 C - zero margin - and contradict the recorded design ruling purely to recolor a gate. A waiver carrying two independent vendor tables is the honest artifact; gate.py records it either way. NOTE: gate waivers are human artifacts by contract (reason + approved required, never agent-invented) - this one is surfaced at H1 and attributed to the owner's 2026-08-16 delegation. \\
P1 & Minimum-load bleed resistor: NOT ruled at P1. Current budget carries 10 mA for it, but whether the part is fitted is deferred to P2/P3 research with a page citation. & The refdesign fragment says LM1117 FIXED variants need no minimum load while AMS1117 specs 3-10 mA, and AMS1117's minimum-load spec is commonly tied to the ADJUSTABLE variant (whose external divider would otherwise supply it). Fitting a bleed the datasheet does not require is padding a block-only board; omitting one it does require lets the output rise out of regulation at no load. Neither error is acceptable and neither is settleable from memory - this is exactly what the coverage/research machinery exists for. Budget stays conservative (10 mA included) so a later YES does not move the thermal numbers. \\
P1 & Two research triggers carried into P2: (a) what AMS1117 Table 1's 'backside ground plane' copper is electrically connected to - until answered only the zero-backside row (1000 mm2 top -\textgreater{} 65 C/W) is usable; (b) the fixed-variant minimum-load requirement. Plus a P3 datasheet-extraction item: actual quiescent current (5 mA assumed; at 10 mA max, Pd = 1.03 W and the target tightens to 63 C/W). & The tab is VOUT, so a via connection from the tab pour to a backside GND plane is impossible - which SUGGESTS Table 1's backside rows are dielectric-only coupling, but that is an inference from the pinout, not a read of the table note. Rows leaning on backside copper stay unusable until a researcher reads the actual page. Designing on an inferred row would put the whole outline on an unverified reading. \\
P2 & Stackup JLC2313\_1.6 (2-layer, 1.6 mm, 1 oz, HASL). 4-layer and 2 oz both rejected. & Architect ruling, accepted: on a tab-on-top part the extra 4L copper is buried 0.5 oz reached through vias and measures \textasciitilde{}20\% worse per mm2 than F.Cu, at 2.5x the board price - more layers do not buy the thing this board needs, which is TOP-SIDE AREA. 2 oz is unquantified for SOT-223 in any source we hold and check\_thermal's model ignores copper weight entirely, so spending it would be faith. 1.6 mm kept because both vendor thermal sweeps were measured on 1/16 in FR-4 - changing thickness would invalidate the curve the board is sized against. \\
P2 & One flat schematic sheet 'main', refdes U1/C1/C2/R1/J1/J2 FIXED (constraints.json keys on them). P4 runs a single schematic agent. No sim bench. & Six parts, one rail, zero signal nets - a hierarchy would be ceremony. Sim: the only interesting analog question is loop stability vs Cout ESR, which is IC-INTERNAL compensation; policy forbids vendored vendor models and a generic datasheet-parameter model card cannot honestly settle a compensation question. A bench that cannot fail truthfully is worse than no bench. \\
P2 & RISK CARRIED INTO P6: board\_edit --outline fit counts courtyards, copper items and rule areas - NOT zones. The entire thermal margin lives in a \textgreater{}=1000 mm2 F.Cu +3V3 ZONE, so a tightly-packed placement would let fit shrink the board below the radiator it needs, silently, with no gate able to see it. Mitigation to apply at P6: reserve the pour area as a rule area (which fit DOES count) before fitting, and re-check credited copper after. & Architect's riskiest-decision call, and it is correct: this is the one failure mode that would produce a board that passes every gate and cooks. It is the bb-buck defect in a new costume - there, a stated size bound the placement; here, an uncounted zone would let the placement give away the size the thermal design earned. \\
P2 & R1 minimum-load bleed resistor is REMOVED from the design. BOM is 5 parts: U1, C1 (Cin), C2 (Cout), J1, J2. & Research finding, independently verified by a fresh second reader off the pages: the 1117 family's minimum-load-current spec belongs to the ADJUSTABLE variant, whose EXTERNAL divider draws it; on a fixed 3.3 V part that divider is internal and Kelvin-connected. Evidence the reader added beyond the researcher's own argument: AMS p1 Ordering Information lists bare 'AMS1117' as its own SOT-223 orderable (the adjustable), AMS p2 conditions the ADJUSTABLE load regulation on 10 mA \textless{}= IOUT \textless{}= 0.8 A while every FIXED row reads 0 \textless{}= IOUT with fixed VOUT rows inheriting an IOUT = 0 mA header, and TI is explicit rather than inferential (LM1117 p6 scopes the minimum-load row to LM1117-ADJ). AMS p6's curve is titled verbatim 'Minimum Operating Current (Adjustable Device)'. No fixed-variant condition anywhere implies a minimum load. Fitting a part the datasheet does not require would pad a block-only board; this deletion is what the coverage machinery bought. \\
P2 & The design MAY credit backside copper thermally. B.Cu GND pour is worth \textasciitilde{}5 C/W at 1000 mm2 and \textasciitilde{}10 C/W at 2500 mm2 under a fixed top pour, WITHOUT any electrical connection to the VOUT tab. Sizing rule unchanged (\textgreater{}=1000 mm2 contiguous F.Cu +3V3 pour, target theta\_JA 65 C/W); the backside credit is taken as MARGIN, not as licence to shrink the top pour. & Verified: AMS1117 p5, in the paragraph directly above Table 1, states verbatim that the heat spreading copper layer 'does not need to be electrically connected to the tab of the device' and names a ground plane 'either inside or on the opposite side of the board' - the second reader confirmed this is NOT scoped to a grounded tab, and re-checked that all six printed Table 1 rows match (55/55/65/80/60/65). So the VOUT tab does not forfeit the backside rows: heat crosses the FR-4 dielectrically. The apparent AMS-vs-TI conflict resolved rather than averaged: TI Fig 9-11 ranks top-Cu \textgreater{} half/half \textgreater{} bottom-Cu because it SPLITS a fixed copper budget, while AMS ADDS backside area under a fixed top pour - different experiments. Top-only columns agree independently (TI 645 mm2 -\textgreater{} 66 C/W, AMS 1000 mm2 -\textgreater{} 65 C/W). Taking it as margin not licence because a 2-layer board gets the bottom GND pour anyway: expect \textasciitilde{}60 C/W and Tj \textasciitilde{}110 C in practice against a 115 C ruling. \\
P2 & The P8 thermal\_area waiver must cite the SOURCES' OWN hedges, not only their numbers: LM1117 section 9 carries TI's boxed note that 'Information in the following applications sections is not part of the TI component specification, and TI does not warrant its accuracy or completeness' - and the copper-area table, both placement figures, the dissipation-vs-ambient curves and the ESR window ALL sit inside section 9. AMS Table 1 carries its own hedge ('a rough guideline... some experimentation will be necessary'). & The waiver's whole argument is that two independent vendor tables beat check\_thermal's clamped model. That argument is only honest if it also carries what those vendors say about their own numbers. Stating the hedges strengthens the waiver rather than weakening it: it shows the evidence was read, not mined. Surfaced by the second reader as a ledger-wide caveat no record had mentioned. \\
P3 & SOT-223 tab = VOUT is CONFIRMED, and the confirmation comes from the LM1117 datasheet, not the AMS one. The extractor was right to flag it as an inference within its own scope; it is not an inference across the workspace's evidence. & The AMS1117 PDF draws 'TAB IS OUTPUT' against the TO-252 front view only; its SOT-223 top view shows leads 1/2/3 and never depicts the tab - so from that document alone, tab identity for SOT-223 is inferred. But the verified record linear-regulator-live-tab-thermal-vias carries TI's LM1117 p3 Figure 6-1 (DCY 4-pin SOT top view) labelled 'Tab is VOUT' with Table 6-1 giving VOUT as SOT-223 pins 2 AND 4, pin 4 being the tab - an explicit, second-reader-verified statement about the SOT-223 package specifically. Recorded so no later stage 'corrects' the extractor's flag by guessing, and so the schematic/footprint stages know the tab is a live VOUT node with cross-vendor backing. KiCad's standard footprint agrees by name: SOT-223-3\_TabPin2. \\
P3 & C1 (10 uF tantalum input capacitor) is RETAINED even though the AMS1117 datasheet specifies no input capacitor anywhere. & The extractor correctly refused to invent one - AMS's Figures 1 and 2 show none and its Stability section is silent. The justification is cross-vendor and already verified: LM1117 p14 (record linear-regulator-1117-output-cap-esr-window, rule key cin\_dielectric=tantalum cin\_uf=10) specifies a 10 uF tantalum at the input for the same topology. Physically it is not optional here either: the input arrives through a screw terminal on bench wire, so the source has metres of lead inductance and no local charge reservoir at VIN. This is support the topology requires at the stated operating point, not added filtering, so it stays in scope at block-only. Recorded because the primary datasheet's silence would otherwise read as 'not required'. \\
P3 & U1's footprint comes from a standard IPC-derived SOT-223 land pattern (KiCad Package\_TO\_SOT\_SMD:SOT-223-3\_TabPin2 or an equivalent built to IPC), NOT from a vendor land drawing - the AMS datasheet has none. & The extractor searched and found only package body/lead dimensions (p7-8), no recommended PCB copper land pattern, and correctly left pitch\_mm/pad\_size\_mm unset rather than deriving pad copper from a package outline (a classic way to get a footprint subtly wrong). A standard IPC SOT-223 land pattern is the right substitute and the librarian must be told the vendor drawing does not exist, so it does not go looking for one or fabricate dimensions. The TAB PAD SIZE is load-bearing on this board - it is the thermal interface to the pour - so the librarian must state where its dimensions came from. \\
P3 & P3 coverage: the escalated first call (floor=proven) re-reported block:B2 as a gap; NOT re-researched. Re-run without the escalation flags: 0 gaps, block:B2 provisional, part:C6186 covered, exit 0. & The gap was produced by the maturity floor alone, not by missing knowledge: B2's five classes are all populated by records this run just researched to VERIFIED, and 'proven' is a floor only bring-up evidence can clear (knowledge.py --prove). build-modes.md is explicit that the post-research re-run drops those flags because re-escalating only re-fires the trigger - firing a second research task at a slot we exhausted four sources on this hour would have spent budget to learn nothing. No coverage-mapper was needed either: the records key the slot directly, so there was no fuzzy remainder to classify. \\
P3 & C2 = Vishay 293D226X9016D2TE3 (C215872), 22 uF / 16 V solid MnO2 tantalum, ESR 800 mohm at 100 kHz. ACCEPTED against the record's window on LCSC's parametric attribute rather than on Vishay's own datasheet page. & The acceptance test (record linear-regulator-1117-output-cap-esr-window) is that the part's ESR sit inside 0.3-22 ohm. 800 mohm is not marginal - 2.7x above the floor and 27x below the ceiling - and the record reserves the bench ring-count settlement for MARGINAL candidates, which this is not. The sourcer surveyed \textasciitilde{}30 orderable 22 uF tantalums with live ESR data and correctly REJECTED the two traps: Vishay's T58 polymer-tantalum at 300 mohm (right at the window edge, and 6.3 V - under the derating floor anyway) and a Panasonic 90 mohm part (ceramic-class ESR, plainly outside). It also overrode blocks.md's tighter 0.3-0.5 ohm first choice under that document's own fallback clause, and said so instead of silently substituting. RESIDUAL, recorded not hidden: the ESR figure is a distributor parametric attribute, not a page read from Vishay's 293D datasheet; confirming it against the manufacturer document is the one open item on this BOM. It does not block P4 - a wrong reading here shows up as ringing at bring-up, which is exactly what the record's ring-count test is for. \\
P3 & J1/J2 = WJ500V-5.08-2P (C8465), 2-position 5.08 mm THT screw terminal, 18 A / 250 V, 0.2-2.5 mm2 wire. Through-hole, hand-soldered after PCBA. & The sourcer confirmed with two targeted queries that JLC's catalog has NO SMT equivalent at this pitch and wire class - so THT is not a preference the mode relaxed, it is the only option, and the hand-solder step is a fact about the catalog rather than a packaging choice. Current rating clears the 500 mA design load by 36x, which costs nothing here because the terminal's size is set by the wire it accepts, not by its current rating. \\
P3 & U1 keeps the LCSC/EasyEDA-pulled SOT-223 footprint, in which the TAB IS ITS OWN PAD NUMBER 4 (KiCad's standard SOT-223-3\_TabPin2 instead reuses pad number 2 for the tab). P4 REQUIREMENT, not a nicety: the schematic must explicitly wire symbol pin 4 to +3V3. & Swapping to the KiCad standard file would also require editing the symbol to drop pin 4 - a two-file structural change on the board's most safety-critical part, for no electrical gain. The pulled symbol/footprint pair is internally consistent and the tab dims (2.340 x 3.600 mm) sanity-check against KiCad's standard (2.0 x 3.8 mm). THE RISK THIS CREATES, and why it is recorded loudly: on this board the tab pad is the thermal interface to the \textasciitilde{}1000 mm2 +3V3 pour. A schematic that wires only pin 2 would leave the tab netless, the pour would never connect to it, and the board would pass ERC while having no heatsink at all - a silent thermal failure that looks clean on every gate that runs before bring-up. Checks: netlist audit at P4, schematic-reviewer, and at P6/P7 confirm the tab pad carries +3V3 and is joined to the pour. \\
P3 & APPROVED a hand edit to three library footprints: redraw the F.CrtYd courtyards of U1, C1 and C2 to the pad-field bounding box + 0.25 mm. J1/J2 untouched (its courtyard already contains both pads). & easyeda2kicad's standing defect: the pulled courtyards trace the PLASTIC BODY, not the copper. Measured overruns - U1 leads \textasciitilde{}2.5 mm and tab \textasciitilde{}2.3 mm outside the courtyard, C1 \textasciitilde{}0.7 mm per side, C2 \textasciitilde{}0.85 mm per side. Two consequences, the second being the one that matters here: (1) P6 placement legality would accept a neighbour sitting on U1's leads; (2)  MEASURES COURTYARDS, so body-sized courtyards bias the EARNED OUTLINE too small - on the one board whose entire thermal margin is copper area. Fixing the library is cheaper and more honest than compensating with a fudge margin at fit time. No fpfix/script covers courtyard geometry (it handles silk gaps only), so this is necessarily a hand edit, made under explicit orchestrator approval and reported with before/after extents. \\
P3 & CORRECTION to the preceding courtyard decision: its consequence (2) was recorded with a clause lost to a shell quoting slip. It should read - board\_edit --outline fit MEASURES COURTYARDS among its content, so body-sized courtyards bias the EARNED OUTLINE too small, on the one board whose entire thermal margin is copper area. & The recorded text read '(2)  MEASURES COURTYARDS' with the script name missing, which obscures the actual mechanism. Correcting rather than leaving an ambiguous audit entry: this is the specific reason the courtyard hand-edit was authorized, and a later reader of state.json needs to see which script does the measuring. \\
P4 & thermal.min\_vias 12 -\textgreater{} 0 in constraints.json. NO thermal vias under U1's tab. B.Cu stays a CONTINUOUS GND plane under the pour. & Schematic-reviewer finding, and it was a real self-contradiction in the constraint set: min\_vias 12 demanded stitching vias in a +3V3 tab pour while B.Cu is GND - a via between those two nets is a SHORT, not a heat path (verified record linear-regulator-live-tab-thermal-vias). The alternative, a bottom-side +3V3 island to stitch into, would punch a hole in the ground return directly under the hottest part of the board to buy copper the dielectric already gives us for free. The record's own prescription for 2 layers over a GND bottom is precisely: large top-side VOUT pour, ground plane continuous beneath, NO vias through it - the far-side copper still works, coupled through the FR-4 (AMS1117 p5), worth about 5 C/W at 1000 mm2. The 12-via number was inherited from exposed-pad-to-GND practice where it is correct; here it is not. constraints\_lint re-run clean. \\
P4 & P6 REQUIREMENT (schematic-reviewer ERROR, deferred to the layer where it lives): silkscreen must legend both screw terminals - polarity (+5V / GND on J1, +3V3 / GND on J2) AND direction (IN / OUT). Scripted silk only (silk\_place.py / place\_edit add\_text), pin-locked. & The reviewer is right and this is in scope: it is about THIS board's own correctness, not about adding an excluded protection feature. The WJ500V footprint's only cues are a 0.2 mm corner dot and two IDENTICAL arrows over both pins, and J1 and J2 are the SAME part - so nothing on the assembled board distinguishes input from output or plus from ground. On a board that by scope ruling has no reverse-polarity protection, a user wiring 5 V backwards puts -5 V on a VIN with no negative abs-max rating and reverse-biases C1, an MnO2 tantalum, with up to 2 A of bench supply behind it: it fails SHORT and can burn. Feeding 5 V into J2 likewise forces VOUT \textgreater{} VIN, also unrated. A legend costs nothing but silkscreen and is the only thing standing between the user and that outcome. Deferred to P6 because silk is a board artifact, not a schematic one - no schematic change is needed. \\
P4 & C2 ESR margin: ACCEPTED as unquantified-at-the-floor, to be settled at bring-up by load-transient ring count (\textless{}= 4 rings), not by the datasheet number. & Reviewer sharpened the earlier residual usefully: LCSC's 0.8 ohm is a 100 kHz CEILING from a parametric field, while the window's 0.3 ohm end is a FLOOR on MINIMUM-over-frequency ESR - so the two numbers are not measured the same way and the margin that actually matters is not quantified by anything we hold. No Vishay 293D datasheet is in research/sources to close it. This does not block: it is not marginal on the face of it (0.8 vs 0.3), and the verified record already prescribes the ring-count test as the settlement for exactly this case. Recorded so bring-up knows it is a real check to run, not a formality. \\
\bottomrule
\end{longtable}
}
\subsection*{Sheet plan}
\subsection*{bb-ldo - hierarchical sheet plan (P2)}
\subsubsection*{One flat sheet. There is no hierarchy on this board.}
Six parts (one contingent), three nets, no signal net, no repeated block. A hierarchical split would add sheet pins and a stitching root to a schematic that fits in one screen. {*}{*}P4 runs a SINGLE schematic agent on the root sheet; there are no child sheets to parallelise.{*}{*}

\begin{quote}\small\ttfamily\raggedright
\textbar{} sheet \textbar{} file \textbar{} blocks \textbar{} parts \textbar{}\\
\textbar{}---\textbar{}---\textbar{}---\textbar{}---\textbar{}\\
\textbar{} `main` (root) \textbar{} `kicad/gen/main.py` -\textgreater{} `kicad/bb-ldo.kicad\_sch` \textbar{} B1 input interface, B2 linear regulator, B3 output interface \textbar{} J1, C1, U1, C2, J2 \textbar{}
\end{quote}
\subsubsection*{Refdes ranges}
Refs must be unique across sheets; with one sheet the whole space is \texttt{main}'s, but the ranges are stated so a later sheet (if this board ever grows) cannot collide.

\begin{quote}\small\ttfamily\raggedright
\textbar{} class \textbar{} range \textbar{} assigned \textbar{}\\
\textbar{}---\textbar{}---\textbar{}---\textbar{}\\
\textbar{} U \textbar{} U1-U9 \textbar{} {*}{*}U1{*}{*} = the linear regulator \textbar{}\\
\textbar{} C \textbar{} C1-C9 \textbar{} {*}{*}C1{*}{*} = 10 uF tantalum Cin, {*}{*}C2{*}{*} = 22 uF solid tantalum Cout \textbar{}\\
\textbar{} R \textbar{} R1-R9 \textbar{} {*}(none fitted - the min-load bleed was RESOLVED OUT at P2; see below){*} \textbar{}\\
\textbar{} J \textbar{} J1-J9 \textbar{} {*}{*}J1{*}{*} = DC input screw terminal, {*}{*}J2{*}{*} = 3V3 output screw terminal \textbar{}\\
\textbar{} \#PWR \textbar{} {*}{*}pwr\_base = 100{*}{*} (\#PWR100-\#PWR199) \textbar{} power symbols \textbar{}
\end{quote}
\textbf{These specific refdes are NOT free choices for P4.} \texttt{constraints.json} already keys on them: \texttt{thermal{[}0{]}.ref = U1}, \texttt{placement.edges} on J1/J2, \texttt{placement.groups.reg} = U1 + C1 + C2, \texttt{placement.sides} on all five certain parts. A different assignment silently unhooks the thermal and placement constraints.

\subsubsection*{Nets - canonical names, fixed here}
\begin{quote}\small\ttfamily\raggedright
\textbar{} net \textbar{} name in the netlist \textbar{} notes \textbar{}\\
\textbar{}---\textbar{}---\textbar{}---\textbar{}\\
\textbar{} input rail \textbar{} `+5V` \textbar{} global power symbol, bare name (no leading `/`) \textbar{}\\
\textbar{} output rail \textbar{} `+3V3` \textbar{} global power symbol; this is ALSO the tab/heatsink net \textbar{}\\
\textbar{} return \textbar{} `GND` \textbar{} global power symbol \textbar{}
\end{quote}
All three are global power symbols, so {*}{*}no hierarchical pins and no local labels exist on this board{*}{*} - nothing acquires a \texttt{/sheet/} prefix. These names are what \texttt{constraints.json} (power, thermal, planes) already declares, and \texttt{netlist\_audit} compares the two: any rename breaks every gate at once.

ERC hint: the \texttt{+5V} and \texttt{GND} nets are driven only by J1's passive connector pins, so both need a PWR\_FLAG (schlib \texttt{power\_flag}) or ERC reports an undriven power net.

\subsubsection*{What each block contributes to the sheet}
\begin{itemize}
\item \textbf{B1} - J1: \texttt{+5V}, \texttt{GND}. Nothing else; no protection element sits behind it
(scope tier).
\item \textbf{B2} - U1 with C1 at its input pin and C2 at its output pin. \textbf{No R1}: the
minimum-load bleed was resolved OUT at P2 on verified, page-cited research -
the 1117 family's minimum-load spec belongs to the ADJUSTABLE variant, whose
EXTERNAL divider draws it, while a fixed 3.3 V part's divider is internal and
Kelvin-connected (AMS p1/p2/p6, LM1117 p6; second-reader verified). The board
carries five parts. {*}{*}Pin numbers come from the P3 datasheet
extraction (\texttt{parts/\textless{}code\textgreater{}.json}), never from this document or from memory{*}{*};
what P2 fixes is the topology: Cin between \texttt{+5V} and \texttt{GND} at the VIN pin,
Cout between \texttt{+3V3} and \texttt{GND} at the VOUT pin, ground pin to \texttt{GND}, and the
SOT-223 tab on \texttt{+3V3} (it is internally VOUT - the footprint's tab pad must
be on the output net, not GND, and not left unconnected).
There is \textbf{no enable pin and no feedback divider} on the fixed variant; if
P3's final part turns out to have an EN pin, it is tied to its always-on
state as a datasheet-required connection (still in scope at \texttt{block-only}),
not brought out as a strap.
\item \textbf{B3} - J2: \texttt{+3V3}, \texttt{GND}.
\end{itemize}
\subsubsection*{decoupling.json (emitted by the generator)}
C1 and C2 are associated with U1's input and output pins. {*}{*}Do NOT tag either with \texttt{"role": "reg\_input"}{*}{*} - that role exists for a SWITCHING regulator's VIN and makes \texttt{check\_decoupling} demand an HF ceramic within 7.5 mm. This block has no switch node, and \texttt{blocks.md} s.3 records why no 0.1 uF ceramic is fitted. Both rails carry a \textgreater{}= 1 uF bulk cap, so \texttt{check\_pdn} has no bulk warning to raise.

\subsubsection*{Placement groups this plan implies (P6)}
\texttt{reg} = U1 (anchor) + C1 + C2, already declared in \texttt{constraints.json{[}"placement"{]}{[}"groups"{]}}. J1 and J2 are edge-pinned to opposite edges. All five certain parts are side-pinned to \texttt{front} (single-sided SMT assembly; the F.Cu pour is the heatsink).

Full architecture narrative (not inlined; see the artifact index): \texttt{architecture/blocks.md}, \texttt{architecture/power\_tree.md}, \texttt{architecture/stackup.md}.
\section{Schematic}
Gate \texttt{erc} (phase P4): \textbf{pass} after 2 attempt(s); last 2026-08-16T23:29:14, failing 0 of 0.
The full schematic PDF follows.
\includepdf[pages=-]{reports/schematic.pdf}
\section{Layout}
\emph{Pending --- produced at P6.}
\section{Verification}
\emph{Pending --- produced at P8.}
\section{DFM and Fabrication}
\emph{Pending --- produced at P9.}
\section{Run Record (Appendix)}
\subsection*{Phase digests}
\paragraph*{P0}
\begin{quote}\small\ttfamily\raggedright
\# P0 Intake - digest\\
~\\
- Mode `learning block-basics:` -\textgreater{} scope {*}{*}block-only{*}{*}, binding {*}{*}canonical{*}{*}\\
  (geometry is an OUTPUT). Recorded; `board\_init` will refuse a fixed outline.\\
- `requirements.md` written; `check\_requirements.py` clean (0 violations).\\
- 9 questions asked in ONE batch, all 9 answered. Safety section 8 closes on\\
  owner answers: bench supply, no battery, no mains-direct, no motor load.\\
- {*}{*}Design point frozen: 0.975 W at Ta = 50 C, still air, continuous.{*}{*}\\
  Vin 4.75-5.25 V -\textgreater{} 1.45 V worst-case headroom, non-LDO admissible.\\
- Output +/-3\% into a static resistive load; datasheet stability caps only.\\
- Tier exclusions accepted (protection/indicators/test-points/config/\\
  second-rail/mechanical). Advisory: hot surface, miswire kills the board.\\
- Owner delegated design judgment + checkpoint ceremony; gates, coverage,\\
  research, DFM and any NEW safety question stay undelegated.\\
- fable unavailable this session -\textgreater{} every fable tier substitutes UP to opus.
\end{quote}
\paragraph*{P1}
\begin{quote}\small\ttfamily\raggedright
\# P1 Research - digest\\
~\\
- Roster: component-scout + reference-design (parallel), then power-architect.\\
  No interface-spec agent: a screw terminal is standards-free.\\
- {*}{*}Lead part AMS1117-3.3 SOT-223, C6186 (Basic, \$0.20).{*}{*} Rejected MCP1825S -\\
  better on paper, but its only theta\_JA (62 C/W) is a JEDEC 4-LAYER number,\\
  unreachable on 2L with no curve to design against. AMS1117 is the only\\
  candidate with a copper-area table on OUR board class (1oz FR-4): 90 C/W min\\
  pad -\textgreater{} 65 at 1000 mm2 -\textgreater{} 55 at 2500 mm2; TI's LM1117 sweep agrees \textasciitilde{}1 C/W.\\
- Design point: Pd 1.00 W, target theta\_JA 65 C/W -\textgreater{} Tj 115 C (121 C stacked\\
  corner). Needs \textgreater{}= 1000 mm2 contiguous F.Cu +3V3 pour, 645 mm2 floor.\\
- Tab = VOUT: heatsink pour is +3V3 and CANNOT be stitched to bottom GND;\\
  B.Cu = GND + a +3V3 island, 12 vias AROUND the pad (via-in-pad wicks solder).\\
- Caps are datasheet requirements: Cout 22 uF SOLID TANTALUM (bare MLCC sits\\
  below the ESR window - the 1117 oscillation trap), Cin 10 uF tantalum.\\
- {*}{*}Gate conflict surfaced now{*}{*}: check\_thermal clamps credited copper at\\
  645 mm2 -\textgreater{} 73.85 C/W floor on 2L, so dt\_c 65 can never pass (verified in\\
  source). Keep 65, waive at P8 with both vendor tables; do NOT relax to 75.\\
- Into P2: what Table 1's backside copper connects to; fixed-variant min-load\\
  requirement; actual Iq (P3). Expected outline \textasciitilde{}35-45 mm/side - not a cap.
\end{quote}
\paragraph*{P2}
\begin{quote}\small\ttfamily\raggedright
\# P2 Architecture + coverage research - digest\\
~\\
- 3 blocks, one rail, zero signal nets. Stackup {*}{*}JLC2313\_1.6{*}{*} (2L, 1 oz).\\
  ONE flat sheet. No sim bench (loop stability is IC-internal compensation).\\
- Coverage at the learning floor (`proven` + `--research-provisional`): 1 slot,\\
  1 GAP - the library had no linear-regulator checklist and no records.\\
- Research `block-linear-regulator-1`: 4 sources (AMS1117 vendor-layout; TI\\
  LM1117 / AN-1028 / SLVA115A cross-vendor), 20 pages read VISUALLY, 6 records\\
  + the topology's first checklist.\\
- {*}{*}3 reader passes, 2 refutations, both the SAME defect{*}{*}: an unsourced\\
  part-selection instruction (first as a PMOS/PNP-scoped ESR-vs-frequency\\
  basis in `rule`, then as "MLCC and polymer are low-ESR by construction" in\\
  prose). A third instance was caught in the CHECKLIST - the artifact that\\
  propagates to future boards - and rewritten to mechanism + acceptance test.\\
  Final 6/6 verified.\\
- Two findings CHANGED the board: {*}{*}R1 deleted{*}{*} (minimum load is the\\
  ADJUSTABLE variant's spec -\textgreater{} 5 parts) and {*}{*}backside copper counts{*}{*} (AMS p5:\\
  the spreading layer need not connect to the tab; \textasciitilde{}5 C/W at 1000 mm2). Taken\\
  as margin, not licence to shrink the top pour.\\
- Coverage re-run: 0 gaps, 1 {*}{*}provisional{*}{*} (owner approval -\textgreater{} covered).\\
- Carried to P8: every TI number here sits inside LM1117 section 9, which TI\\
  disclaims; AMS Table 1 self-describes as "a rough guideline".
\end{quote}
\paragraph*{P3}
\begin{quote}\small\ttfamily\raggedright
\# P3 Parts + Library - digest\\
~\\
- BOM = 4 distinct parts / 5 placements, \textasciitilde{}\$1.05/board at qty 5: U1\\
  AMS1117-3.3 SOT-223 (C6186, Basic), C1 10 uF/16 V solid Ta (C7171, Basic),\\
  C2 22 uF/16 V solid Ta (C215872), J1=J2 WJ500V-5.08-2P (C8465).\\
- {*}{*}C2 was the judgment call{*}{*}: acceptance test is ESR INSIDE 0.3-22 ohm (too\\
  LITTLE ESR oscillates this part). Chose 800 mohm; rejected a 300 mohm\\
  polymer (window edge, 6.3 V) and a 90 mohm part (ceramic-class). Residual:\\
  LCSC parametric attribute, not a Vishay page. Not marginal, not blocking.\\
- J1/J2 THT is forced, not preferred - JLC has no SMT terminal at this pitch.\\
- C6186.json clean; it correctly REFUSED to invent a land pattern (the AMS\\
  datasheet has none, only package body dims).\\
- {*}{*}Two findings that could have killed the board silently:{*}{*}\\
  1. The pulled footprint gives the TAB its own pad 4 (KiCad's standard reuses\\
     pad 2). The schematic MUST wire pin 4 to +3V3, or the tab is netless, the\\
     pour never connects, and ERC passes on a board with no heatsink.\\
  2. easyeda2kicad drew courtyards around the PLASTIC BODY (U1 leads 2.57 mm\\
     outside). Hand-edited under approval to pad bbox + 0.25 mm on U1/C1/C2.\\
     Matters because `--outline fit` measures courtyards.\\
- fp\_verify U1 pass post-edit; C1/C2 polarity confirmed pin 1 = positive.
\end{quote}
\paragraph*{P4}
\begin{quote}\small\ttfamily\raggedright
\# P4 Schematic - digest\\
~\\
- One flat sheet, 5 parts, 3 nets, 12/12 pins connected. Generator is the\\
  source (`kicad/gen/main.py`); the .kicad\_sch is build output and rebuilds\\
  byte-identical from a clean delete.\\
- {*}{*}Gate erc: 0 errors / 0 warnings{*}{*} (commit 55d7875). netlist\_audit exit 0\\
  with one expected `power\_no\_consumers` on +3V3 (its only consumer is the\\
  off-board load through passive J2).\\
- `U1 \{1:GND, 2:+3V3, 3:+5V, 4:+3V3\}` - the TAB (pin 4) is on +3V3, read back\\
  from the netlist, not from source.\\
- {*}{*}The catch of the run{*}{*}: the readability redraw initially left the tab run\\
  unnamed. Pins 2 and 4 are one node INSIDE the package but two nodes on the\\
  sheet, so the run exported as Net-(C2-Pad1) with +3V3 holding only pin 2 -\\
  the netless-tab thermal failure, arriving through a COSMETIC edit. ERC said\\
  0/0 (an unnamed 3-pin net is legal). Only `netlist\_audit --compare` saw it.\\
- Reviewer (fresh context): 1 error, 5 warnings. Circuit itself clean - tab\\
  net, both tantalum anodes, C2's compensation role, abs-max and dropout all\\
  verified and NOT reported.\\
  - ERROR: no +/GND or IN/OUT legend on J1/J2 (same part, identical arrows).\\
    Deferred to P6 as a silk requirement - board layer, not schematic.\\
  - Fixed here: constraints min\_vias 12 -\textgreater{} 0 (a +3V3 pour cannot stitch to a\\
    GND plane); schematic redrawn to real wires; text collisions removed.\\
  - Accepted: C2's ESR margin is unquantified at the floor -\textgreater{} bring-up ring\\
    count (\textless{}= 4) settles it.\\
- 3 library defects repaired idempotently in `kicad/gen/lib\_fixups.py` (C2 pin\\
  angles, C1 non-polarized graphic, numeral pin names on all 2-pin symbols).\\
  MUST be re-run after any lib\_pull refresh or they return.
\end{quote}
\subsection*{Gate attempt history}
{\small
\begin{longtable}{llll}
\toprule
\textbf{Timestamp} & \textbf{Gate} & \textbf{Status} & \textbf{Failing/Total} \\
\midrule
\endhead
2026-08-16T22:42:46 & erc & pass & 0/0 \\
2026-08-16T23:29:14 & erc & pass & 0/0 \\
\bottomrule
\end{longtable}
}
\subsection*{Issues}
\emph{(no entries)}
\subsection*{Human checkpoints}
{\small
\begin{longtable}{lllp{2.6cm}p{6.4cm}}
\toprule
\textbf{Id} & \textbf{Phase} & \textbf{Status} & \textbf{When} & \textbf{Note} \\
\midrule
\endhead
1 & P2 & approved & 2026-08-16T19:51:53 & Presented as a digest under the owner's 2026-08-16 delegation ('make decisions yourself based on what will give the best learning for the basic system'); no separate owner sign-off was solicited or given, and the run proceeded. Items that still need a GENUINE owner ruling and are NOT covered by this: promotion/approval of the 6 verified records + the linear-regulator checklist into the shared library, the P8 thermal\_area waiver (gate waivers are human artifacts by contract), and any order at P10. \\
2 & P4 & not reached &  &  \\
3 & P6 & not reached &  &  \\
4 & P8 & not reached &  &  \\
5 & P10 & not reached &  &  \\
\bottomrule
\end{longtable}
}
\section{Artifact Index}
{\small
\begin{longtable}{p{9cm}rp{4cm}}
\toprule
\textbf{File} & \textbf{Size} & \textbf{Note} \\
\midrule
\endhead
\texttt{state.json} & 46,774 B &  \\
\texttt{kicad/bb-ldo.kicad\_sch} & 37,327 B &  \\
\texttt{brief/brief.md} & 608 B &  \\
\texttt{requirements.md} & 13,344 B &  \\
\texttt{architecture/blocks.md} & 16,013 B &  \\
\texttt{architecture/constraints.json} & 5,928 B &  \\
\texttt{architecture/power\_tree.md} & 8,341 B &  \\
\texttt{architecture/sheets.md} & 4,242 B &  \\
\texttt{architecture/stackup.md} & 6,862 B &  \\
\texttt{reports/gate-erc.json} & 10,065 B &  \\
\texttt{reports/review-schematic.json} & 3,920 B &  \\
\texttt{reports/review-schematic.md} & 11,453 B &  \\
\texttt{reports/schematic.pdf} & 47,921 B &  \\
\bottomrule
\end{longtable}
}
\end{document}
